\documentclass{article}
\usepackage{amsmath,amssymb}
\usepackage[margin=1in]{geometry}
\allowdisplaybreaks
\begin{document}

\section*{Model M: common-information profit sharing, final $\alpha$ formulation}

\begin{equation}
\begin{aligned}
\max_{\alpha}\quad &
\frac{(E(P_s)/2-f-\alpha)\sigma(D)\beta}
{\left(E(D)+\dfrac{\alpha}{\sqrt{\beta-\alpha^2}}\sigma(D)\right)(\beta-\alpha^2)^{3/2}}
\left(\alpha E(D)-\sigma(D)\sqrt{\beta-\alpha^2}+\frac{E(P_sD)}{2}\right)\\
\text{subject to}\quad &
\alpha\in\left[\frac{E(P_s)}{2}-w_u,\frac{E(P_s)}{2}-w_l\right].
\end{aligned}
\tag{21}
\end{equation}
The displayed objective is the paper's expression for $(1-\gamma)\Pi$, representing the retailer's profit $(1-\gamma)\Pi(w(\gamma))$. The maximization and interval are specified in the prose immediately following equation (21), physical PDF page 22.

\paragraph{Printed endpoint names and derived identification.}
The authors do not separately define the printed endpoint names $w_u$ and $w_l$. Physical PDF page 22 states $w(Q)\in[f,w_0]$ and $\alpha=E(P_s)/2-w$. The image of that price interval under this order-reversing affine map is $[E(P_s)/2-w_0,E(P_s)/2-f]$. Comparing it with the printed interval in (21) gives the derived endpoint identifications $w_u=w_0$ and $w_l=f$. These are source-supported derived mappings, not a claim that the authors explicitly printed those aliases. The quantity $w_0$ remains the globally determined no-sharing response defined below, not the threshold $w^{UB}$ or an arbitrary replacement bound. The distinct quantity endpoints $Q_l,Q_u$, printed objective and recovery formulas, and their endpoint and undefined-value qualifications remain as stated below.

\subsection*{Fixed inputs, information, and decision roles}
Demand $D\geq0$ and selling price $P_s\geq0$ are exogenous random quantities. Both parties have the same prescribed first and second moments. The supplier's fixed unit cost satisfies $f\geq0$. The retailer chooses the profit share in stage 1, the supplier chooses the unit wholesale price $w\geq0$ in stage 2, and the retailer chooses capacity $Q\geq0$ in stage 3. The wholesale payment is linear:
\[
c(Q)=wQ.
\]
Unused capacity has zero salvage value. Before taking the worst case over distributions, the decentralized profit and the parties' shares are
\begin{equation}
\Pi=E_{P_s,D}\!\left[P_s\min(Q,D)\right]-wQ.
\tag{2}
\end{equation}
\[
\Pi_R:=(1-\gamma)\Pi,\qquad
\Pi_S:=\gamma\Pi+(w-f)Q.
\]
Here $\gamma$ is the supplier's share and $1-\gamma$ is the retailer's share; $\alpha$ in (21) is the transformed wholesale-price variable, not a profit share.

The given moments matrix and its distributional ambiguity set are
\[
\boldsymbol{\Sigma}=
\begin{pmatrix}
E(P_s^2)&E(P_sD)&E(P_s)\\
E(P_sD)&E(D^2)&E(D)\\
E(P_s)&E(D)&1
\end{pmatrix},
\]
\[
\Phi(\boldsymbol{\Sigma}_0)=
\left\{(P_s,D)\;;\;P_s,D\geq0,\space{}
P_s,D\text{ have moments matrix }\boldsymbol{\Sigma}=\boldsymbol{\Sigma}_0\right\}.
\]
The matrix $\boldsymbol{\Sigma}_0$ is fixed data. All five moments in these formulas take their values from its entries. As stated on page 10, a moments matrix of nonnegative $P_s,D$ is positive semidefinite and has nonnegative entries:
\[
\boldsymbol{\Sigma}\succeq0,\qquad \boldsymbol{\Sigma}\geq0.
\]
The latter inequality is entrywise. The notation $\sigma(D)$ denotes the standard deviation of demand.

\subsection*{Retailer response and the selected closed-form replacement}
The original response and value to which the final formulation refers are
\begin{equation}
Q(w):=\operatorname{argmax}_{Q\geq0}
\left\{-wQ+\inf_{(P_s,D)\in\Phi(\boldsymbol{\Sigma}_0)}
E_{P_s,D}\!\left(P_s\min(Q,D)\right)\right\},
\tag{3}
\end{equation}
\begin{equation}
\Pi(w):=\max_{Q\geq0}
\left\{-wQ+\inf_{(P_s,D)\in\Phi(\boldsymbol{\Sigma}_0)}
E_{P_s,D}\!\left(P_s\min(Q,D)\right)\right\}.
\tag{4}
\end{equation}
These specify the meaning of the response functions; the final $\alpha$ formulation uses the authors' closed-form response below, not an additional semidefinite optimization problem.

Definition 1 (page 15) gives
\[
\alpha(w)=\frac{E(P_s)}{2}-w,\qquad
\beta=\frac{E(P_s^2)}{4}.
\]
Theorem 3 and Definition 2 (page 16) give the threshold and the mandatory input premise
\[
w^{UB}=\frac12\left[
E(P_s)+
\frac{E(P_sD)E(D)-\sigma(D)\sqrt{E(P_s^2)E(D^2)-E(P_sD)^2}}
{E(D^2)}
\right],\qquad w^{UB}>f.
\]
Under the printed condition
\begin{equation}
w\leq\frac12\left[
E(P_s)+
\frac{E(P_sD)E(D)-\sigma(D)\sqrt{E(P_s^2)E(D^2)-E(P_sD)^2}}
{E(D^2)}
\right],
\tag{17}
\end{equation}
the selected retailer response is
\begin{equation}
Q(w)=Q^*=E(D)+\frac{\alpha(w)}{\sqrt{\beta-\alpha(w)^2}}\sigma(D).
\tag{18}
\end{equation}
The associated value, using Theorem 3's replacement of $Q_r^*$ by $Q^*$ in equation (16), is
\begin{equation}
\begin{aligned}
\Pi(w)=\Pi_{wst}
&=-wQ(w)+\inf_{(P_s,D)\in\Phi(\boldsymbol{\Sigma}_0)}
E_{P_s,D}\!\left(P_s\min(D,Q(w))\right)\\
&=\alpha(w)E(D)-\sigma(D)\sqrt{\beta-\alpha(w)^2}
+\frac12E(P_sD),\qquad w\leq w^{UB}.
\end{aligned}
\tag{16}
\end{equation}
Theorem 3 states that otherwise $Q^*=0$. Its proof (pages 47--48) states
\[
\Pi(w)=0\quad\text{for }w\geq w^{UB},\qquad
Q^*=0\quad\text{for }w>w^{UB}.
\]
In particular, at $w=w^{UB}$ the printed closed-form response is retained and the profit is zero. Section 5 considers wholesale prices inducing positive retailer orders. Appendix B, page 41, states that the expression for the positive response requires $\beta>\alpha^2$; this is not used here to replace or tighten the printed interval in (21). No extra convention for zero denominators or degenerate moments is introduced.

\subsection*{Supplier response retained in the final formulation}
Section 5 restricts the supplier's response to $f\leq w\leq w^{UB}$ and defines
\[
\Pi_1(w)=(w-f)Q(w),\qquad
\Pi_S(w,\gamma)=\Pi_1(w)+\gamma\Pi(w),\quad w\leq w^{UB}.
\]
\begin{equation}
w(\gamma):=\operatorname{argmax}_{f\leq w\leq w^{UB}}
\left\{\Pi_1(w)+\gamma\Pi(w)\right\}.
\tag{19}
\end{equation}
This is the supplier response used by the selected version, replacing the earlier unrestricted response over $w\geq0$. For $w>w^{UB}$, Section 5 prints $\Pi_S=0$.

The final response analysis treats $\gamma\in[0,1]$. In particular,
\[
w(1)=f,\qquad w(\gamma)>f\quad\text{for all }\gamma\in[0,1).
\]
At $\gamma=1$, the retailer is indifferent because its profit share is zero. Endnote 3 specifies the Strong Stackelberg Equilibrium convention: among multiple optimal retailer decisions, the retailer chooses $Q$ to maximize the supplier's payoff.

\paragraph{Scope of the share range.}
The starting equation (5) prints only $\gamma\geq0$, whereas Sections 5.1--5.2 and Theorem 4 use shares in $[0,1]$. The selected formulation is the final equation (21), with its printed $\alpha$ interval and the endpoint response rules below; the earlier $\gamma$ maximization is not imposed as another objective.

\subsection*{Response inversion, endpoint definitions, and recovery}
The inversion printed on page 21 is
\[
\alpha(Q)=
\begin{cases}
\displaystyle\frac{\sqrt{\beta}}{\sqrt{1+\left(\dfrac{\sigma(D)}{Q-E(D)}\right)^2}},&Q\geq E(D),\\[6pt]
\displaystyle-\frac{\sqrt{\beta}}{\sqrt{1+\left(\dfrac{\sigma(D)}{Q-E(D)}\right)^2}},&Q<E(D),
\end{cases}
\]
\[
w(Q)=\frac{E(P_s)}2-\alpha(Q)=
\begin{cases}
\displaystyle\frac{E(P_s)}2-\frac{\sqrt{\beta}}{\sqrt{1+\left(\dfrac{\sigma(D)}{Q-E(D)}\right)^2}},&Q\geq E(D),\\[6pt]
\displaystyle\frac{E(P_s)}2+\frac{\sqrt{\beta}}{\sqrt{1+\left(\dfrac{\sigma(D)}{Q-E(D)}\right)^2}},&Q<E(D).
\end{cases}
\]
Definition 3 and the surrounding text give
\[
w_0=w(\gamma=0),\qquad f<w_0\leq w^{UB},\qquad
Q_u=Q(f),\qquad Q_l=Q(w_0),
\]
\[
w(Q)\in[f,w_0],\qquad Q\in[Q_l,Q_u].
\]
For the interior response, the paper uses the first-order condition
\[
\frac{\partial\Pi_S}{\partial w}(w^*,\gamma)
=\frac{d\Pi_1}{dw}(w^*)+\gamma\frac{d\Pi}{dw}(w^*)=0
\]
and states
\begin{equation}
\gamma(Q)=-\frac{\dfrac{d\Pi_1}{dw}(w(Q))}
{\dfrac{d\Pi}{dw}(w(Q))}.
\tag{20}
\end{equation}
The derivative definitions used in Appendix B, pages 48--49, are
\begin{align}
\frac{dQ}{dw}(w)&=-\frac{\beta\sigma(D)}{(\beta-\alpha^2)^{3/2}}<0,
&&f\leq w\leq w^{UB},\tag{48}\\
\frac{d\Pi}{dw}(w)&=-\left(E(D)+\sigma(D)\frac{\alpha}{\sqrt{\beta-\alpha^2}}\right)=-Q(w),\tag{49}\\
\frac{d\Pi_1}{dw}(w)&=Q(w)+(w-f)\frac{dQ}{dw}(w),\tag{50}\\
\frac{\partial\Pi_S}{\partial w}(w,\gamma)&=(1-\gamma)Q(w)+(w-f)\frac{dQ}{dw}(w),\tag{51}
\end{align}
where $\alpha=\alpha(w)$ in these derivative formulas.

Theorem 4 states that the outcome is $Q\in[Q_l,Q_u]$ when the retailer chooses
\[
\gamma=1-
\frac{(E(P_s)/2-f-\alpha(Q))\sigma(D)\beta}
{(\beta-\alpha(Q)^2)^{3/2}Q},\qquad Q\in(Q_l,Q_u),
\]
with the separate endpoint choices
\[
Q=Q_l:\ \gamma=0,\qquad Q=Q_u:\ \gamma=1.
\]
The text also states $\gamma(Q)=0$ if and only if $w^*=w_0$, and $\gamma(Q)=1$ if and only if $w^*=f$.

For an optimizer $\alpha^*$ of (21), the printed recovery equation is
\begin{equation}
\gamma^*=1-
\frac{(E(P_s)/2-f-\alpha^*)\sigma(D)\beta}
{\left(E(D)+\dfrac{\alpha^*}{\sqrt{\beta-(\alpha^*)^2}}\sigma(D)\right)
(\beta-(\alpha^*)^2)^{3/2}}.
\tag{22}
\end{equation}

\paragraph{Boundary and sufficiency qualifications retained from the source.}
The branch $Q\geq E(D)$ in the printed inversion includes $Q=E(D)$ although its expression divides by $Q-E(D)$; no value at that equality is supplied, and none is added here. Page 21 explicitly calls the first-order condition necessary but potentially insufficient, then asserts that the share in (20) induces the desired supplier response. Theorem 4 retains that assertion and supplies the endpoint choices separately. The objective (21) and recovery (22) are printed over the closed $\alpha$ interval without separate endpoint qualifications, while the explicit interior share formula in Theorem 4 is scoped to $Q\in(Q_l,Q_u)$. These scopes are preserved rather than repaired. No unavailable external supplement or erratum has been used; any correction from such a source remains unverified.

\end{document}
